\section[Khảo sát]{Khảo sát}

\subsection{2.1 Nghiệp vụ}

\begin{frame}{Olist kết nối người bán nhỏ với các sàn lớn: 3.095 seller bán 99.441 đơn trong 2 năm}
\muc{2.1 · Nghiệp vụ}
\begin{itemize}
  \item Dữ liệu Brazil, 09/2016–10/2018; hai bộ Kaggle
\end{itemize}
\khungchen[0.3\textheight]{CHÈN HÌNH / BẢNG: ảnh/logo Olist}
\canxacnhan
\note{[S13] Nhóm tìm hiểu Olist từ nguồn công khai; .}
\end{frame}

\begin{frame}{Mô hình kinh doanh hai phía: seller được thu hút, người mua mua qua sàn}
\muc{2.1 · Nghiệp vụ}
\begin{columns}[T]\begin{column}{0.4\textwidth}
\begin{itemize}
  \item Business Model Canvas 9 ô (bản nháp)
\end{itemize}
\end{column}\begin{column}{0.58\textwidth}
\begin{center}\anh[0.74\textheight]{3-2-3_business_model_canvas.png}\end{center}
\end{column}\end{columns}
\canxacnhan
\note{[S14] Ô doanh thu là suy luận, dữ liệu không có doanh thu của Olist.}
\end{frame}

\begin{frame}{Giá trị đề xuất: một hợp đồng, nhiều sàn; dữ liệu cho thấy phí vận chuyển bằng 16,6\% giá hàng}
\muc{2.1 · Nghiệp vụ}
\begin{columns}[T]\begin{column}{0.4\textwidth}
\begin{itemize}
  \item Giá trị cho seller và cho người mua; chi phí chính là vận chuyển và hậu cần
\end{itemize}
\end{column}\begin{column}{0.58\textwidth}
\begin{center}\anh[0.74\textheight]{3-2-3_business_model_canvas.png}\end{center}
\end{column}\end{columns}

\note{[S15] GMV 13,49 triệu BRL, phí vận chuyển 2,24 triệu BRL.}
\end{frame}

\begin{frame}{Bài toán: hiểu hiệu suất đơn hàng và hiệu quả thu hút seller từ hai hệ thống rời nhau}
\muc{2.1 · Nghiệp vụ}
\begin{columns}[T]\begin{column}{0.4\textwidth}
\begin{itemize}
  \item Hai hệ thống chỉ gặp nhau ở \texttt{seller\_\allowbreak{}id}
\end{itemize}
\end{column}\begin{column}{0.58\textwidth}
\begin{center}\anh[0.74\textheight]{3-2-4_erd_oltp.png}\end{center}
\end{column}\end{columns}

\note{[S16] Giá trị của kho là tích hợp để trả lời câu hỏi xuyên quy trình.}
\end{frame}

\begin{frame}{Sáu nhóm câu hỏi phân tích bao trùm bán hàng, vận hành, khách hàng và thu hút seller}
\muc{2.1 · Nghiệp vụ}
\begin{columns}[T]\begin{column}{0.4\textwidth}
\begin{itemize}
  \item Bán hàng $\cdot$ xử lý đơn $\cdot$ thanh toán $\cdot$ trải nghiệm khách $\cdot$ seller $\cdot$ thu hút seller
\end{itemize}
\end{column}\begin{column}{0.58\textwidth}
\begin{center}\anh[0.74\textheight]{3-2-6_cay_yeu_cau_phan_tich.png}\end{center}
\end{column}\end{columns}

\note{[S17] Mỗi nhóm trở thành một hoặc nhiều dashboard.}
\end{frame}

\begin{frame}{Kho phục vụ hai quy trình nghiệp vụ: xử lý đơn và thu hút seller}
\muc{2.1 · Nghiệp vụ}
\begin{columns}[T]\begin{column}{0.4\textwidth}
\begin{itemize}
  \item Chung chiều seller và ngày
\end{itemize}
\end{column}\begin{column}{0.58\textwidth}
\begin{center}\anh[0.74\textheight]{3-3-4_constellation.png}\end{center}
\end{column}\end{columns}

\note{[S18] Chiều seller là "chiếc cầu" nối hai quy trình.}
\end{frame}

\begin{frame}{Đơn hàng đi qua bốn mốc thời gian và có thể bị hủy hoặc không khả dụng}
\muc{2.1 · Nghiệp vụ}
\begin{columns}[T]\begin{column}{0.4\textwidth}
\begin{itemize}
  \item Đặt $\rightarrow$ duyệt $\rightarrow$ bàn giao vận chuyển $\rightarrow$ giao $\rightarrow$ đánh giá; 1.234 đơn Failed (1,24\%)
\end{itemize}
\end{column}\begin{column}{0.58\textwidth}
\begin{center}\anh[0.74\textheight]{3-2-1_bpmn_don_hang.png}\end{center}
\end{column}\end{columns}

\note{[S19] BPMN nháp; [nhóm vẽ lại bằng draw.io].}
\end{frame}

\begin{frame}{Mỗi mốc của quy trình là một cột thời gian mà kho biến thành chỉ số}
\muc{2.1 · Nghiệp vụ}
\begin{itemize}
  \item \texttt{approval\_\allowbreak{}hours} 10,3 giờ $\cdot$ \texttt{handling\_\allowbreak{}days} 2,9 ngày $\cdot$ \texttt{transit\_\allowbreak{}days} 9,3 ngày $\cdot$ \texttt{lead\_\allowbreak{}days} 12,6 ngày
\end{itemize}
\begin{center}\adjustbox{max width=\linewidth,max height=0.5\textheight}{\scriptsize\renewcommand{\arraystretch}{1.15}\begin{tabular}{lll}\toprule
\textbf{Chặng} & \textbf{Chỉ số của kho} & \textbf{Trung bình} \\\midrule
Mua $\rightarrow$ duyệt & \texttt{approval\_hours} & 10,3 giờ \\
Duyệt $\rightarrow$ bàn giao & \texttt{handling\_days} & 2,9 ngày \\
Bàn giao $\rightarrow$ giao & \texttt{transit\_days} & 9,3 ngày \\
Mua $\rightarrow$ giao & \texttt{lead\_days} & 12,6 ngày \\
Ngày hứa giao & \texttt{promised\_days} & 24,4 ngày \\
\bottomrule\end{tabular}}\end{center}

\note{[S20] Chặng vận chuyển dài nhất; chặng seller bàn giao gần 3 ngày.}
\end{frame}

\begin{frame}{Kho thu thập dữ liệu ở bảy điểm của quy trình đơn hàng}
\muc{2.1 · Nghiệp vụ}
\begin{columns}[T]\begin{column}{0.4\textwidth}
\begin{itemize}
  \item Ô vàng trên BPMN: đặt, duyệt, bàn giao, giao, đánh giá, hủy, thanh toán
\end{itemize}
\end{column}\begin{column}{0.58\textwidth}
\begin{center}\anh[0.74\textheight]{3-2-1_bpmn_don_hang.png}\end{center}
\end{column}\end{columns}

\note{[S21] Đây là chỗ hệ thống phân tích "bám" vào quy trình.}
\end{frame}

\begin{frame}{Quy trình thu hút seller đi từ MQL đến seller có đơn, trung bình mất 100 ngày}
\muc{2.1 · Nghiệp vụ}
\begin{columns}[T]\begin{column}{0.4\textwidth}
\begin{itemize}
  \item Liên hệ $\rightarrow$ chốt (TB 48,5 ngày) $\rightarrow$ đơn đầu (TB 51,9 ngày)
\end{itemize}
\end{column}\begin{column}{0.58\textwidth}
\begin{center}\anh[0.74\textheight]{3-2-2_bpmn_thu_hut_seller.png}\end{center}
\end{column}\end{columns}

\note{[S22] Hai khoảng thời gian này là hai chỉ số của fact lead funnel.}
\end{frame}

\begin{frame}{Phễu thu hút: 8.000 MQL $\rightarrow$ 842 deal (10,5\%) $\rightarrow$ 379 seller có đơn bán được (45,0\% số deal)}
\muc{2.1 · Nghiệp vụ}
\begin{columns}[T]\begin{column}{0.4\textwidth}
\begin{itemize}
  \item 462 deal đã chốt nhưng chưa bán gì
\end{itemize}
\end{column}\begin{column}{0.58\textwidth}
\begin{center}\anh[0.74\textheight]{3-3-1_funnel_seller_overlap.png}\end{center}
\end{column}\end{columns}

\note{[S23] Đây là kết quả phân tích, không phải lỗi dữ liệu.}
\end{frame}

\begin{frame}{Cây nghiệp vụ gom hoạt động của Olist thành năm nhóm}
\muc{2.1 · Nghiệp vụ}
\begin{columns}[T]\begin{column}{0.4\textwidth}
\begin{itemize}
  \item Thu hút seller $\cdot$ quản lý đơn hàng $\cdot$ chăm sóc khách $\cdot$ quản lý seller $\cdot$ quản trị dữ liệu
\end{itemize}
\end{column}\begin{column}{0.58\textwidth}
\begin{center}\anh[0.74\textheight]{3-2-5_cay_nghiep_vu.png}\end{center}
\end{column}\end{columns}

\note{[S24] Kho chủ yếu phục vụ nhóm quản lý đơn hàng và thu hút seller.}
\end{frame}

\begin{frame}{Mười một điểm phân tích A1--A11 gắn mỗi bước của quy trình với dashboard và KPI}
\muc{2.1 · Nghiệp vụ}
\begin{itemize}
  \item A1--A6 trên quy trình đơn: đặt hàng và thanh toán, duyệt, bàn giao, giao hàng, đánh giá, hủy
  \item A7--A10 trên quy trình seller: MQL, SDR và SR, chốt deal, đơn đầu tiên; A11: giám sát nạp dữ liệu
  \item Ví dụ A4 (giao hàng) $\rightarrow$ DB2, DB3 $\rightarrow$ Avg lead days, On-time rate, Promised vs actual $\rightarrow$ bang, vùng, khoảng cách
\end{itemize}

\note{[S122] Mỗi điểm phân tích truy ngược được từ dashboard về một bước nghiệp vụ cụ thể (bảng trong báo cáo mục 2.1).}
\end{frame}

\begin{frame}{Mỗi khối Business Model Canvas dẫn tới một câu hỏi quản trị, một chỉ số và một dashboard}
\muc{2.1 · Nghiệp vụ}
\begin{itemize}
  \item Giá trị đề xuất $\rightarrow$ giao đúng hạn và đủ rẻ? $\rightarrow$ On-time rate, Freight ratio 16,6\% $\rightarrow$ DB2, DB3
  \item Kênh $\rightarrow$ kênh nào đưa về seller chốt được? $\rightarrow$ MQL to won rate $\rightarrow$ DB4
  \item Dòng doanh thu và cấu trúc chi phí: dữ liệu chỉ có giá hàng và phí vận chuyển nên dùng GMV và tỷ lệ phí làm chỉ số thay thế
\end{itemize}

\note{[S123] Chín khối Canvas, mỗi khối một câu hỏi; hai khối thiếu dữ liệu được nêu rõ.}
\end{frame}

\subsection{2.2 Yêu cầu phân tích}

\begin{frame}{Cây yêu cầu phân tích đi từ sáu chủ điểm đến các câu hỏi cụ thể}
\muc{2.2 · Yêu cầu phân tích}
\begin{columns}[T]\begin{column}{0.4\textwidth}
\begin{itemize}
  \item Chủ điểm $\rightarrow$ câu hỏi $\rightarrow$ chiều cần quan tâm
\end{itemize}
\end{column}\begin{column}{0.58\textwidth}
\begin{center}\anh[0.74\textheight]{3-2-6_cay_yeu_cau_phan_tich.png}\end{center}
\end{column}\end{columns}

\note{[S25] Cây này là cơ sở thiết kế fact, dimension và dashboard.}
\end{frame}

\begin{frame}{Bán hàng: GMV cần tăng trưởng MoM, YoY và quy luật 80/20}
\muc{2.2 · Yêu cầu phân tích}
\begin{itemize}
  \item GMV 13,49 triệu BRL; AOV 137,42 BRL
\end{itemize}
\khungchen[0.3\textheight]{CHÈN HÌNH / BẢNG: bảng chủ điểm bán hàng (3.2.2)}

\note{[S26] GMV không tính đơn hủy để phản ánh doanh thu thật.}
\end{frame}

\begin{frame}{Xử lý đơn: không chỉ hỏi "đúng hạn không" mà hỏi "hứa so với thực tế chênh bao nhiêu"}
\muc{2.2 · Yêu cầu phân tích}
\begin{itemize}
  \item Đúng hạn 93,2\% nhưng hứa 24,4 ngày, thực tế 12,6 ngày
\end{itemize}
\khungchen[0.3\textheight]{CHÈN HÌNH / BẢNG: bảng chủ điểm xử lý đơn}

\note{[S27] Vì ngày dự kiến rất rộng nên nhóm thêm KPI ngày hứa so với ngày giao thực tế.}
\end{frame}

\begin{frame}{Trải nghiệm khách: giao trễ là yếu tố kéo điểm review xuống mạnh nhất}
\muc{2.2 · Yêu cầu phân tích}
\begin{columns}[T]\begin{column}{0.4\textwidth}
\begin{itemize}
  \item Điểm 4,29 khi giao sớm, 1,70 khi trễ trên 7 ngày
\end{itemize}
\end{column}\begin{column}{0.58\textwidth}
\begin{center}\anh[0.74\textheight]{3-3-1_delay_vs_review.png}\end{center}
\end{column}\end{columns}

\note{[S28] Đây là phát hiện trung tâm của DB2 (trang trải nghiệm khách).}
\end{frame}

\begin{frame}{Thu hút seller: câu hỏi là kênh nào chuyển đổi tốt và seller từ funnel bán ra sao}
\muc{2.2 · Yêu cầu phân tích}
\begin{columns}[T]\begin{column}{0.4\textwidth}
\begin{itemize}
  \item Tỷ lệ chốt theo kênh từ 2,7\% đến 16,7\%
\end{itemize}
\end{column}\begin{column}{0.58\textwidth}
\begin{center}\anh[0.74\textheight]{3-3-1_funnel_win_rate_by_origin.png}\end{center}
\end{column}\end{columns}

\note{[S29] Luôn nêu cỡ mẫu vì nhiều kênh chỉ có vài chục seller.}
\end{frame}

\begin{frame}{Tám điểm hiện trạng của hệ thống nguồn dẫn tới tám yêu cầu cải tiến cho kho dữ liệu}
\muc{2.2 · Yêu cầu phân tích}
\begin{itemize}
  \item Hai hệ thống rời nhau, chỉ gặp qua seller\_id $\rightarrow$ chiều dùng chung dim\_seller
  \item Trạng thái đơn không có lịch sử mốc $\rightarrow$ snapshot tích luỹ 5 mốc
  \item 45 loại lỗi dữ liệu $\rightarrow$ ODS gắn cờ, ghi dòng bị loại kèm lý do
  \item Ngày dự kiến giao trung bình 24,4 ngày so với thực tế 12,6 ngày $\rightarrow$ KPI ngày hứa so với thực tế
\end{itemize}

\note{[S124] Bảng hiện trạng và yêu cầu cải tiến ở báo cáo mục 2.2.}
\end{frame}

\begin{frame}{Sáu nhóm người dùng, mỗi nhóm có câu hỏi, dashboard và tần suất xem riêng}
\muc{2.2 · Yêu cầu phân tích}
\begin{itemize}
  \item Ban điều hành $\rightarrow$ DB1 (hằng tuần, hằng tháng); quản lý kinh doanh $\rightarrow$ DB1
  \item Vận hành và logistics $\rightarrow$ DB2, DB3 (hằng ngày đến hằng tuần); chăm sóc khách hàng $\rightarrow$ DB2
  \item Quản lý seller và đội sales $\rightarrow$ DB4; quản trị dữ liệu dùng bảng giám sát ETL ở mục 3.3 (không còn dashboard)
\end{itemize}

\note{[S125] Đối tượng sử dụng, câu hỏi, dashboard và tần suất xem.}
\end{frame}

\subsection{2.3 KPI và dashboard}

\begin{frame}{Bảng KPI chính gồm 19 chỉ số, mỗi chỉ số có một định nghĩa DAX duy nhất}
\muc{2.3 · KPI}
\begin{itemize}
  \item GMV, Orders, AOV, on-time, lead time, review, repeat, MQL$\rightarrow$seller...
\end{itemize}
\begin{columns}[T]\begin{column}{0.49\textwidth}\adjustbox{max width=\linewidth}{\scriptsize\begin{tabular}{lrr}\toprule \textbf{KPI} & \textbf{Toàn bộ dữ liệu} & \textbf{Cửa sổ 2017-01 → 2018-08} \\ \midrule
GMV & 13.494.400,74 & 13.449.529,68 \\
Orders & 98.199 & 97.905 \\
AOV & 137,42 & 137,37 \\
Freight ratio & 16,6\% & 16,6\% \\
Orders placed & 99.441 & 99.092 \\
Cancel rate & 1,24\% & 1,19\% \\
On-time rate & 93,2\% & 93,2\% \\
Avg lead days & 12,56 & 12,54 \\
Avg promised days & 24,37 & 24,28 \\
Promised vs actual (days) & 11,81 & 11,75 \\
\bottomrule\end{tabular}}\end{column}\begin{column}{0.49\textwidth}\adjustbox{max width=\linewidth}{\scriptsize\begin{tabular}{lrr}\toprule \textbf{KPI} & \textbf{Toàn bộ dữ liệu} & \textbf{Cửa sổ 2017-01 → 2018-08} \\ \midrule
Avg review & 4,09 & 4,09 \\
\% negative & 14,7\% & 14,6\% \\
Repeat customer rate & 3,12\% & 3,11\% \\
MQLs (theo tháng liên hệ đầu) & 8.000 & 8.000 \\
Won deals & 842 & 842 \\
MQL to won rate & 10,5\% & 10,5\% \\
Avg days to close & 48,50 & 48,50 \\
Sellers with first sale & 379 & 379 \\
Won to first sale rate & 45,0\% & 45,0\% \\
\bottomrule\end{tabular}}\end{column}\end{columns}

\note{[S30] Giá trị trong bảng tính bằng SQL và dùng để đối chiếu với Power BI.}
\end{frame}

\begin{frame}{Bốn dashboard có bốn cây phân tích từ mục tiêu đến chỉ số và chiều}
\muc{2.3 · KPI}
\begin{itemize}
  \item Bốn cây phân tích DB1--DB4
\end{itemize}
\begin{columns}[c]\begin{column}{0.24\textwidth}\includegraphics[width=\linewidth,height=0.4\textheight,keepaspectratio]{2-3_cay_phan_tich_DB1.png}\end{column}\begin{column}{0.24\textwidth}\includegraphics[width=\linewidth,height=0.4\textheight,keepaspectratio]{2-3_cay_phan_tich_DB2.png}\end{column}\begin{column}{0.24\textwidth}\includegraphics[width=\linewidth,height=0.4\textheight,keepaspectratio]{2-3_cay_phan_tich_DB3.png}\end{column}\begin{column}{0.24\textwidth}\includegraphics[width=\linewidth,height=0.4\textheight,keepaspectratio]{2-3_cay_phan_tich_DB4.png}\end{column}\end{columns}

\note{[S31] Ví dụ DB2: mục tiêu giám sát giao hàng và trải nghiệm khách, chỉ số on-time, promised vs actual, review, chiều vùng và nhóm trễ.}
\end{frame}

\begin{frame}{Năm giới hạn của dữ liệu làm một KPI có hai con số: GMV 13,49 so với 13,59 triệu BRL}
\muc{2.3 · KPI}
\begin{itemize}
  \item GMV loại đơn Failed 13.494.400,74; tổng giá hàng mọi đơn 13.591.643,70 BRL
  \item Orders 98.199 so với Orders placed 99.441 (chênh 1.242 đơn); 380 so với 379 seller có đơn
  \item 2016 chỉ 329 đơn, 2018-09 chỉ 16 đơn; 547 đơn có từ 2 review; điểm đúng hạn từng là 4,21 rồi sửa thành 4,29
  \item Xử lý: chọn một định nghĩa, ghi cạnh con số, sửa ở ETL
\end{itemize}

\note{[S126] Bảng giới hạn dữ liệu ảnh hưởng tới KPI ở báo cáo mục 2.3.}
\end{frame}

\subsection{2.4 Dữ liệu nguồn}

\begin{frame}{Phạm vi: dữ liệu 01/2017–08/2018 chiếm 99,6\% số đơn nên được dùng cho xu hướng}
\muc{2.4 · Dữ liệu nguồn}
\begin{columns}[T]\begin{column}{0.4\textwidth}
\begin{itemize}
  \item 2016 chỉ 329 đơn; 09–10/2018 chỉ 20 đơn
\end{itemize}
\end{column}\begin{column}{0.58\textwidth}
\begin{center}\anh[0.74\textheight]{3-3-1_orders_per_month.png}\end{center}
\end{column}\end{columns}

\note{[S32] Bộ lọc InWindow trong Power BI cố định cửa sổ này cho DB1–DB4.}
\end{frame}

\begin{frame}{Hai hệ thống nguồn có 11 bảng, 121 MB và 1.559.764 dòng}
\muc{2.4 · Dữ liệu nguồn}
\begin{itemize}
  \item oltp 9 bảng 1.550.922 dòng; crm 2 bảng 8.842 dòng
\end{itemize}
\begin{center}\adjustbox{max width=\linewidth,max height=0.5\textheight}{\scriptsize\begin{tabular}{llrrr}
\toprule
\textbf{Bảng nguồn} & \textbf{File CSV} & \textbf{Số dòng} & \textbf{Số cột} & \textbf{MB} \\
\midrule
oltp.product\_category\_name\_translation & product\_category\_name\_translation.csv & 71 & 2 & 0,00 \\
oltp.products & olist\_products\_dataset.csv & 32.951 & 9 & 2,27 \\
oltp.sellers & olist\_sellers\_dataset.csv & 3.095 & 4 & 0,17 \\
oltp.geolocation & olist\_geolocation\_dataset.csv & 1.000.163 & 6 & 58,44 \\
oltp.customers & olist\_customers\_dataset.csv & 99.441 & 5 & 8,62 \\
oltp.orders & olist\_orders\_dataset.csv & 99.441 & 8 & 16,84 \\
oltp.order\_items & olist\_order\_items\_dataset.csv & 112.650 & 7 & 14,72 \\
oltp.order\_payments & olist\_order\_payments\_dataset.csv & 103.886 & 5 & 5,51 \\
oltp.order\_reviews & olist\_order\_reviews\_dataset.csv & 99.224 & 7 & 13,78 \\
crm.mql & olist\_marketing\_qualified\_leads\_dataset.csv & 8.000 & 4 & 0,67 \\
crm.closed\_deals & olist\_closed\_deals\_dataset.csv & 842 & 14 & 0,16 \\
Tổng (11 bảng) &  & 1.559.764 &  & 121,18 \\
\bottomrule
\end{tabular}}\end{center}

\note{[S33] Vượt tiêu chí kích thước dữ liệu của môn (trên 100 MB).}
\end{frame}

\begin{frame}{Bảng toạ độ chiếm 64\% số dòng nhưng chỉ đại diện 19.015 mã bưu chính}
\muc{2.4 · Dữ liệu nguồn}
\begin{itemize}
  \item 1.000.163 dòng geolocation
\end{itemize}
\khungchen[0.3\textheight]{CHÈN HÌNH / BẢNG: bảng nguồn (làm nổi dòng geolocation)}

\note{[S34] Đây là bảng cần xử lý nhiều nhất ở ODS.}
\end{frame}

\begin{frame}{Luồng dữ liệu và tần suất cập nhật: nạp ban đầu 45.430 đơn rồi 11 đợt hằng tháng}
\muc{2.4 · Dữ liệu nguồn}
\begin{columns}[T]\begin{column}{0.4\textwidth}
\begin{itemize}
  \item Tổng 286.144 bản ghi mới qua 11 đợt
\end{itemize}
\end{column}\begin{column}{0.58\textwidth}
\begin{center}\anh[0.74\textheight]{3-4-1_pipeline.png}\end{center}
\end{column}\end{columns}

\note{[S37] Dữ liệu Kaggle là ảnh tĩnh nên nhóm mô phỏng cập nhật.}
\end{frame}

\begin{frame}{Trung bình 163 đơn mỗi ngày, đỉnh 1.176 đơn ngày Black Friday: lô hằng ngày là đủ}
\muc{2.4 · Dữ liệu nguồn}
\begin{columns}[T]\begin{column}{0.4\textwidth}
\begin{itemize}
  \item Cửa sổ 01/2017--08/2018: 163 đơn, 185 dòng hàng, 13 MQL, 1,3 deal mỗi ngày
  \item Bảng tần suất đề xuất công nghệ: lô hằng ngày bằng procedure PostgreSQL; gần thời gian thực thì CDC + Kafka (chương 5)
\end{itemize}
\end{column}\begin{column}{0.58\textwidth}
\begin{center}\anh[0.74\textheight]{2-4_tan_suat_cap_nhat.png}\end{center}
\end{column}\end{columns}

\note{[S127] Biểu đồ tần suất cập nhật theo thời gian và bảng công nghệ xử lý đề xuất.}
\end{frame}

\begin{frame}{Cây phân loại dữ liệu chia 11 bảng thành 4 nhóm: chủ, giao dịch, tương tác, tham chiếu}
\muc{2.4 · Dữ liệu nguồn}
\begin{columns}[T]\begin{column}{0.4\textwidth}
\begin{itemize}
  \item Mỗi bảng kèm hai cột tiêu biểu và giá trị thật
  \item Phân loại thêm theo kiểu dữ liệu, độ nhạy cảm (định danh, vị trí) và tính biến động
\end{itemize}
\end{column}\begin{column}{0.58\textwidth}
\begin{center}\anh[0.74\textheight]{2-4_data_taxonomy.png}\end{center}
\end{column}\end{columns}

\note{[S128] Data taxonomy của 11 bảng nguồn.}
\end{frame}

\begin{frame}{Data flow hiện trạng: mũi tên ghi thông tin được truyền, hai hệ thống nguồn tách rời}
\muc{2.4 · Dữ liệu nguồn}
\begin{columns}[T]\begin{column}{0.4\textwidth}
\begin{itemize}
  \item Thông tin đơn, trạng thái giao, thanh toán, điểm đánh giá, lead
  \item Marketing funnel (SDR, SR) tách khỏi hệ thống bán hàng
\end{itemize}
\end{column}\begin{column}{0.58\textwidth}
\begin{center}\anh[0.74\textheight]{2-4_data_flow_hien_trang.png}\end{center}
\end{column}\end{columns}

\note{[S129] Hiện trạng: luồng thông tin giữa các hệ thống CNTT (bản nháp, nhóm vẽ lại).}
\end{frame}

\begin{frame}{Data flow đề xuất: nguồn, staging, ODS, kho, mart và BI, có nhật ký ETL giám sát}
\muc{2.4 · Dữ liệu nguồn}
\begin{columns}[T]\begin{column}{0.4\textwidth}
\begin{itemize}
  \item Bảng lớp, thành phần, công nghệ, tần suất, ghi chú ở báo cáo
  \item Điều phối bằng một lệnh Python, biến đổi bằng PL/pgSQL
\end{itemize}
\end{column}\begin{column}{0.58\textwidth}
\begin{center}\anh[0.74\textheight]{2-4_data_flow_de_xuat.png}\end{center}
\end{column}\end{columns}

\note{[S130] Đề xuất: luồng nguồn đến BI.}
\end{frame}

\subsection{2.5 ERD OLTP}

\begin{frame}{ERD hệ thống bán hàng: chín bảng nối quanh đơn hàng}
\muc{2.5 · ERD OLTP}
\begin{columns}[T]\begin{column}{0.4\textwidth}
\begin{itemize}
  \item orders $\rightarrow$ items, payments, reviews; customers, sellers, products
\end{itemize}
\end{column}\begin{column}{0.58\textwidth}
\begin{center}\anh[0.74\textheight]{3-2-4_erd_oltp.png}\end{center}
\end{column}\end{columns}

\note{[S35] \texttt{customer\_\allowbreak{}id} được sinh riêng cho từng đơn nên cần \texttt{customer\_\allowbreak{}unique\_\allowbreak{}id} để nhận ra khách thật.}
\end{frame}

\begin{frame}{Các cột được gạch chân trên ERD là cột đi sang kho, kèm ghi chú mapping}
\muc{2.5 · ERD OLTP}
\begin{itemize}
  \item Ví dụ: \texttt{order\_\allowbreak{}purchase\_\allowbreak{}timestamp} $\rightarrow$ \texttt{purchase\_\allowbreak{}date\_\allowbreak{}key}, \texttt{price} $\rightarrow$ \texttt{price}
\end{itemize}
\khungchen[0.3\textheight]{CHÈN HÌNH / BẢNG: ERD đã gạch chân (vẽ lại)}

\note{[S36] Bảng mapping đầy đủ 108 dòng ODS $\rightarrow$ kho nằm ở phụ lục.}
\end{frame}

\begin{frame}{ERD tô màu theo nhóm dữ liệu: 65 trong 71 cột nguồn được gạch chân vì đi sang OLAP}
\muc{2.5 · ERD OLTP}
\begin{columns}[T]\begin{column}{0.4\textwidth}
\begin{itemize}
  \item Gạch chân suy ra từ bảng mapping nguồn $\rightarrow$ ODS $\rightarrow$ kho
  \item 6 cột không dùng: geolocation.id, 3 thuộc tính sản phẩm, 2 thuộc tính khai báo của deal
\end{itemize}
\end{column}\begin{column}{0.58\textwidth}
\begin{center}\anh[0.74\textheight]{2-5_erd_oltp_gach_chan.png}\end{center}
\end{column}\end{columns}

\note{[S131] ERD nháp tô màu và gạch chân cột dùng sang OLAP.}
\end{frame}

\begin{frame}{Luồng OLTP sang OLAP cho thấy bảng nguồn nào nuôi dimension nào và fact nào}
\muc{2.5 · ERD OLTP}
\begin{columns}[T]\begin{column}{0.4\textwidth}
\begin{itemize}
  \item Bảng tổng hợp mapping theo bảng nguồn ở báo cáo mục 2.5
  \item Mapping đầy đủ từng cột ở Phụ lục C
\end{itemize}
\end{column}\begin{column}{0.58\textwidth}
\begin{center}\anh[0.74\textheight]{2-5_data_flow_oltp_olap.png}\end{center}
\end{column}\end{columns}

\note{[S132] Sơ đồ luồng dữ liệu OLTP sang OLAP.}
\end{frame}

